\documentclass[12pt,a4paper]{article}
\usepackage[utf8]{inputenc}
\usepackage[english]{babel}
\usepackage{amsmath}
\usepackage{amssymb}
\usepackage{graphicx}
\usepackage{hyperref}
\usepackage{listings}
\usepackage{xcolor}
\usepackage{geometry}
\usepackage{booktabs}
\usepackage{algorithm}
\usepackage{algorithmic}

\geometry{margin=2.5cm}

\lstset{
    basicstyle=\ttfamily\small,
    breaklines=true,
    frame=single,
    backgroundcolor=\color{gray!10},
    keywordstyle=\color{blue}\bfseries,
    commentstyle=\color{green!60!black},
    stringstyle=\color{red}
}

\title{\textbf{CVFormer: Transformer Autoencoder for Extracting Collective Variables from Molecular Dynamics Trajectories}\\[0.5cm]
\large Complete Technical Documentation}
\author{Research Project}
\date{\today}

\begin{document}

\maketitle
\tableofcontents
\newpage

\section{Introduction}

\subsection{Context and Motivation}
CVFormer is a Transformer-based system for automatic extraction of collective variables (CVs) from molecular dynamics (MD) trajectories. Collective variables are reduced-dimensional coordinates that describe the relevant conformational states of a biomolecular system, essential for enhanced sampling techniques such as metadynamics.

The fundamental challenge addressed by CVFormer is: \textit{given a protein trajectory with hundreds of conformational degrees of freedom (dihedral angles), how can we automatically identify a small set (typically 2-3) of meaningful coordinates that capture the essential conformational dynamics?}

Traditional approaches require manual selection of CVs based on chemical intuition or extensive trial-and-error. CVFormer automates this process through:
\begin{enumerate}
    \item \textbf{Learning from data}: the model observes thousands of conformational snapshots from MD simulations
    \item \textbf{Self-supervised compression}: it learns to compress all dihedral information into a low-dimensional latent space while preserving reconstruction accuracy
    \item \textbf{Interpretability}: attention mechanisms reveal which residues are most important for defining conformational states
    \item \textbf{Direct integration}: the learned CVs can be immediately used in production MD simulations via PLUMED
\end{enumerate}

\subsection{The Core Idea: Why This Works}

The key insight is that protein conformational changes are \textit{not random}. Despite having $N$ residues each with 2 dihedral angles ($\phi$, $\psi$), most conformational variation occurs in a much lower-dimensional manifold. CVFormer discovers this manifold through:

\begin{itemize}
    \item \textbf{Bottleneck autoencoder}: forcing information through a narrow latent space (2-3 dimensions) ensures that only the most essential conformational features are retained
    \item \textbf{Attention mechanism}: automatically identifies which residues contribute most to conformational variability
    \item \textbf{Periodic representation}: using $(\sin, \cos)$ pairs respects the circular nature of dihedral angles
\end{itemize}

In practice, this means CVFormer \textit{learns} what an expert would manually determine: that perhaps only 3-5 key residues in a hinge region or binding site determine most of the protein's conformational behavior.

\subsection{Main Features}
\begin{itemize}
    \item \textbf{Transformer Architecture}: uses self-attention mechanisms to capture long-range dependencies between residues
    \item \textbf{Attention Pooling}: aggregates residue information through attention weights, allowing identification of the most important residues
    \item \textbf{Unit circle projection}: normalization of $(\sin, \cos)$ pairs to ensure numerical stability
    \item \textbf{Mask handling}: automatic support for residues without valid $\phi/\psi$ angles (e.g., termini)
    \item \textbf{Contrastive Learning}: optional SimCLR-style contrastive loss to improve representation robustness
\end{itemize}

\section{System Architecture}

\subsection{File Structure}
The project is organized in a modular structure:

\begin{lstlisting}[language=bash]
cvformer/
├── pkgs/                           # Core modules
│   ├── model.py                    # Transformer architecture
│   ├── train.py                    # Training loop and metrics
│   └── utils.py                    # Preprocessing and dataset
├── testit.py                       # Main training script
├── environment.yml                 # Conda dependencies
├── data/                           # Sample data
├── cvformer-test/                  # Tests and experiments
│   ├── analyze/                    # Analysis notebooks
│   └── test*/                      # Test results
└── readme.md                       # Documentation
\end{lstlisting}

\subsection{Transformer Autoencoder Model: The Complete Pipeline}

\subsubsection{Philosophy: Why an Autoencoder?}

An autoencoder architecture is ideal for CV discovery because:
\begin{enumerate}
    \item \textbf{Self-supervised}: requires only trajectory data, no manual labels
    \item \textbf{Information bottleneck}: the narrow latent space forces the model to discard noise and retain only essential features
    \item \textbf{Reconstruction guarantee}: if the model can reconstruct all dihedrals from 2-3 latent values, those values must capture the conformational state
\end{enumerate}

The workflow is: \textit{dihedral angles} $\rightarrow$ \textbf{encoder} $\rightarrow$ \textit{latent CVs} $\rightarrow$ \textbf{decoder} $\rightarrow$ \textit{reconstructed dihedrals}

During training, the model learns to make the latent space as informative as possible. After training, we discard the decoder and use only the encoder to compute CVs on-the-fly in MD simulations.

\subsubsection{Input Representation: Respecting Periodicity}

The dihedral angles $\phi$ and $\psi$ of each residue are \textit{periodic} ($-\pi$ is the same as $+\pi$). Using raw angle values would create artificial discontinuities. Instead, we project onto the unit circle:
\begin{equation}
\mathbf{x}_i = [\sin \phi_i, \cos \phi_i, \sin \psi_i, \cos \psi_i] \in \mathbb{R}^4
\end{equation}

This representation:
\begin{itemize}
    \item Preserves continuity: $\phi = -\pi$ and $\phi = +\pi$ map to the same point
    \item Ensures bounded inputs: all values are in $[-1, 1]$
    \item Maintains angle recovery: $\phi = \text{atan2}(\sin \phi, \cos \phi)$
\end{itemize}

The model input is therefore a sequence of $N$ tokens (residues):
\begin{equation}
\mathbf{X} \in \mathbb{R}^{B \times N \times 4}
\end{equation}
where $B$ is the batch size and $N$ is the number of residues. Each token represents one residue's backbone conformation.

\subsubsection{Encoder: From Residues to Collective Variables}

The encoder transforms a sequence of $N$ residue conformations into $d_{\text{latent}}$ collective variables (typically 2-3). It proceeds through five stages:

\paragraph{Stage 1: Input Embedding}
The 4-dimensional $(\sin, \cos)$ representation is projected into a higher-dimensional space ($d_{\text{model}} = 64$-$128$):
\begin{equation}
\mathbf{H}^{(0)} = \mathbf{X} \mathbf{W}_{\text{in}} + \mathbf{b}_{\text{in}} \quad \in \mathbb{R}^{B \times N \times d_{\text{model}}}
\end{equation}

\textit{Why higher dimension?} The Transformer needs rich representations to compute attention. The projection expands the feature space so the model can learn complex relationships between residues.

\paragraph{Stage 2: Positional Encoding}
Since Transformers have no inherent notion of sequence order, we add position information:
\begin{equation}
PE_{(pos, 2i)} = \sin\left(\frac{pos}{10000^{2i/d_{\text{model}}}}\right), \quad
PE_{(pos, 2i+1)} = \cos\left(\frac{pos}{10000^{2i/d_{\text{model}}}}\right)
\end{equation}

This encoding is \textit{learnable-free} and allows the model to distinguish residue 5 from residue 50, even though both might have similar $\phi/\psi$ values. Position matters in proteins: the same angle in a loop vs. a helix has different structural implications.

\paragraph{Stage 3: Self-Attention Layers}
The core of the encoder is a stack of $L_{\text{enc}}$ Transformer layers (typically 3-4). Each layer performs:
\begin{align}
\text{Attention}(\mathbf{Q}, \mathbf{K}, \mathbf{V}) &= \text{softmax}\left(\frac{\mathbf{Q}\mathbf{K}^T}{\sqrt{d_k}}\right)\mathbf{V} \\
\mathbf{H}^{(\ell+1)} &= \text{TransformerLayer}(\mathbf{H}^{(\ell)})
\end{align}

\textit{What does self-attention do?} For each residue, it computes a weighted combination of \textit{all other residues}, allowing the model to capture:
\begin{itemize}
    \item \textbf{Local dependencies}: neighboring residues in secondary structures
    \item \textbf{Long-range interactions}: residues far apart in sequence but close in 3D space
    \item \textbf{Cooperative motions}: groups of residues that move together
\end{itemize}

Each layer refines these representations. Early layers capture local patterns; deeper layers capture global conformational features.

\paragraph{Stage 4: Attention Pooling (The Critical Step)}
We now have $N$ refined token representations $\mathbf{h}_i$. To compress to a fixed-size latent vector, we use \textit{learned attention pooling}:
\begin{align}
s_i &= \mathbf{v}^T \tanh(\mathbf{W}_a \mathbf{h}_i) \quad \text{(score each residue)} \\
\alpha_i &= \frac{\exp(s_i)}{\sum_{j \in \text{valid}} \exp(s_j)} \quad \text{(normalize to weights)} \\
\mathbf{z}_{\text{pool}} &= \sum_{i=1}^{N} \alpha_i \mathbf{h}_i \quad \text{(weighted average)}
\end{align}

\textit{This is where interpretability emerges.} The weights $\alpha_i$ indicate how much each residue contributes to the latent representation. Residues with high $\alpha$ are the most important for defining conformational states. For example:
\begin{itemize}
    \item In a domain-motion protein: hinge residues get high weights
    \item In a binding pocket: active-site residues dominate
    \item In an allosteric protein: key coupling residues stand out
\end{itemize}

\paragraph{Stage 5: Latent Projection}
Finally, the pooled vector is mapped to the target latent dimension:
\begin{equation}
\mathbf{z} = \text{MLP}(\mathbf{z}_{\text{pool}}) \in \mathbb{R}^{d_{\text{latent}}}
\end{equation}

The MLP consists of: Linear $\rightarrow$ LayerNorm $\rightarrow$ GELU $\rightarrow$ Dropout $\rightarrow$ Linear. This produces the final collective variables $\mathbf{z}$ that will be used for biasing simulations.

\subsubsection{Decoder: Reconstructing the Full Conformation}

The decoder takes the compressed latent vector $\mathbf{z}$ and reconstructs all $N$ residue dihedrals. This acts as a \textit{quality check}: if reconstruction is accurate, the latent space is sufficiently informative.

\paragraph{Architecture Design Choice}
Unlike simple autoencoders that directly expand $\mathbf{z}$ to output dimension, CVFormer uses a proper \textit{Transformer decoder} with cross-attention. This allows the model to:
\begin{itemize}
    \item Flexibly allocate information across residues based on learned queries
    \item Handle variable-length proteins (via learned queries)
    \item Leverage the same powerful attention mechanisms used in encoder
\end{itemize}

\paragraph{Stage 1: Latent Expansion to Memory}
The latent vector is expanded into $M$ ``memory tokens'' (typically 4):
\begin{equation}
\mathbf{M} = \text{MLP}(\mathbf{z}) \rightarrow \text{reshape} \in \mathbb{R}^{M \times d_{\text{model}}}
\end{equation}

\textit{Why multiple memory tokens?} A single vector would create an information bottleneck. $M$ tokens provide multiple ``channels'' through which different aspects of the conformation can flow (e.g., one for secondary structure, one for overall compactness, etc.).

\paragraph{Stage 2: Learned Residue Queries}
For each of the $N$ residues, the decoder has a \textit{learned query vector} $\mathbf{q}_i$. These are trainable parameters initialized randomly and optimized during training:
\begin{equation}
\mathbf{Q} = \mathbf{Q}_{\text{learned}} + PE \in \mathbb{R}^{N \times d_{\text{model}}}
\end{equation}

Each query $\mathbf{q}_i$ learns to ``ask'' the memory for information relevant to residue $i$. Positional encoding is added so queries know their position in the sequence.

\paragraph{Stage 3: Cross-Attention Decoder}
The decoder runs $L_{\text{dec}}$ Transformer decoder layers (typically 3-4), where each layer performs:
\begin{enumerate}
    \item \textbf{Self-attention among queries}: queries communicate with each other (allowing inter-residue context)
    \item \textbf{Cross-attention to memory}: each query attends to all memory tokens to extract relevant information
    \item \textbf{Feedforward}: non-linear processing of the attended information
\end{enumerate}

Mathematically:
\begin{equation}
\mathbf{H}_{\text{dec}} = \text{TransformerDecoder}(\mathbf{Q}, \mathbf{M})
\end{equation}

\textit{What is the decoder learning?} It learns to distribute the compressed information in $\mathbf{z}$ (via memory $\mathbf{M}$) back to individual residues in a spatially coherent way.

\paragraph{Stage 4: Output Projection}
Each residue representation is projected to 4 raw values:
\begin{equation}
\hat{\mathbf{X}}_{\text{raw}} = \mathbf{H}_{\text{dec}} \mathbf{W}_{\text{out}} \in \mathbb{R}^{N \times 4}
\end{equation}

\paragraph{Stage 5: Unit Circle Normalization}
To ensure outputs are valid $(\sin, \cos)$ pairs (i.e., lie on the unit circle), we normalize each pair:
\begin{equation}
\hat{\mathbf{x}}_{i,\phi} = \frac{\hat{\mathbf{x}}_{i,\phi}^{\text{raw}}}{\|\hat{\mathbf{x}}_{i,\phi}^{\text{raw}}\|_2}, \quad
\hat{\mathbf{x}}_{i,\psi} = \frac{\hat{\mathbf{x}}_{i,\psi}^{\text{raw}}}{\|\hat{\mathbf{x}}_{i,\psi}^{\text{raw}}\|_2}
\end{equation}

This guarantees $\sin^2 + \cos^2 = 1$ exactly, respecting the geometric constraint. Without this, the model might output $(0.9, 0.8)$ which doesn't correspond to any valid angle.

\subsection{Loss Function: Training the Autoencoder}

\subsubsection{Reconstruction Loss: The Primary Objective}
The main training objective is to minimize reconstruction error:
\begin{equation}
\mathcal{L}_{\text{rec}} = \frac{1}{|\mathcal{M}| \cdot B \cdot 4} \sum_{b=1}^{B} \sum_{i \in \mathcal{M}} \|\hat{\mathbf{x}}_{b,i} - \mathbf{x}_{b,i}\|_2^2
\end{equation}
where $\mathcal{M}$ is the set of valid residues (mask = True), $B$ is batch size.

\textit{Why MSE on $(\sin, \cos)$?} This loss:
\begin{itemize}
    \item \textbf{Respects periodicity}: errors near $\pm \pi$ are handled correctly
    \item \textbf{Is differentiable}: enables gradient-based optimization
    \item \textbf{Weights all angles equally}: no bias toward specific angle ranges
\end{itemize}

\textit{How does this force good CVs?} The only way the model can achieve low reconstruction loss is by encoding \textit{all essential conformational information} in $\mathbf{z}$. Since $\mathbf{z}$ is low-dimensional (2-3D), the model must learn to:
\begin{itemize}
    \item Discard noise and local fluctuations
    \item Retain only the degrees of freedom that actually vary across conformations
    \item Implicitly identify the ``slow'' collective motions that distinguish conformational states
\end{itemize}

\subsubsection{Contrastive Loss: Improving Robustness (Optional)}
To make the latent space more robust to noise and improve generalization, CVFormer optionally implements a SimCLR-style contrastive loss.

\paragraph{The Idea}
Two slightly perturbed versions of the same conformation should map to \textit{similar} latent vectors, while truly different conformations should map \textit{far apart}. This encourages:
\begin{itemize}
    \item \textbf{Invariance}: latent space ignores small thermal fluctuations
    \item \textbf{Smoothness}: nearby conformations have nearby latent coordinates
    \item \textbf{Discrimination}: distinct conformational states are well-separated
\end{itemize}

\paragraph{Implementation}
\begin{enumerate}
    \item \textbf{Data Augmentation}: For each conformation $\mathbf{x}$ in a batch, create a noisy version $\mathbf{x}'$:
    \begin{equation}
    \mathbf{x}' = \text{normalize}(\mathbf{x} + \epsilon), \quad \epsilon \sim \mathcal{N}(0, \sigma^2)
    \end{equation}
    where normalization projects back to unit circle. Typical $\sigma = 0.05$.

    \item \textbf{Encode both versions}:
    \begin{equation}
    \mathbf{z} = \text{encoder}(\mathbf{x}), \quad \mathbf{z}' = \text{encoder}(\mathbf{x}')
    \end{equation}

    \item \textbf{NT-Xent Loss}: For each positive pair $(\mathbf{z}_i, \mathbf{z}_i')$, maximize their similarity relative to all other samples in the batch:
    \begin{equation}
    \mathcal{L}_{\text{NT-Xent}} = -\frac{1}{2B} \sum_{i=1}^{B} \left[\log \frac{\exp(\text{sim}(\mathbf{z}_i, \mathbf{z}_i')/\tau)}{\sum_{j \neq i} \exp(\text{sim}(\mathbf{z}_i, \mathbf{z}_j)/\tau)}\right]
    \end{equation}
    where $\text{sim}(\mathbf{u}, \mathbf{v}) = \frac{\mathbf{u}^T \mathbf{v}}{\|\mathbf{u}\|\|\mathbf{v}\|}$ is cosine similarity, and $\tau$ is a temperature parameter (typically 0.1).
\end{enumerate}

\paragraph{Combined Loss}
The total loss is:
\begin{equation}
\mathcal{L} = \mathcal{L}_{\text{rec}} + \lambda_c \mathcal{L}_{\text{NT-Xent}}
\end{equation}
with $\lambda_c \approx 0.1$ (recommended). The reconstruction loss ensures informativeness; the contrastive loss ensures robustness.

\subsection{Evaluation Metrics}

\subsubsection{Angular Mean Absolute Error}
To evaluate reconstruction quality in physically meaningful terms:

\begin{align}
\phi_{\text{pred}} &= \text{atan2}(\sin \phi_{\text{pred}}, \cos \phi_{\text{pred}}) \\
\text{MAE}_\phi &= \frac{1}{|\mathcal{M}| \cdot B} \sum_{b,i \in \mathcal{M}} |\Delta_{\text{circ}}(\phi_{\text{pred}}, \phi_{\text{true}})|
\end{align}

where $\Delta_{\text{circ}}(a, b) = \text{atan2}(\sin(a-b), \cos(a-b))$ is the circular difference in $[-\pi, \pi]$.

\section{Data Preprocessing: From Trajectory to Training Data}

\subsection{The Challenge: Aligning Phi and Psi Angles}

When MDTraj computes dihedral angles, it returns separate arrays for $\phi$ and $\psi$. However, these arrays have \textit{different lengths} and \textit{different residue orderings}:
\begin{itemize}
    \item $\phi$ is undefined for the first residue (no $C_{i-1}$ atom)
    \item $\psi$ is undefined for the last residue (no $N_{i+1}$ atom)
    \item The indices don't directly correspond to the same residue
\end{itemize}

For training, we need a consistent representation where \textit{token $i$ always represents the same residue for both $\phi$ and $\psi$}. The \texttt{compute\_aligned\_phi\_psi()} function solves this.

\subsection{MDTraj Dihedral Conventions}

MDTraj defines dihedrals using 4-atom quadruples:

\begin{itemize}
    \item \textbf{Phi ($\phi$)}: $[C_{i-1}, N_i, CA_i, C_i]$

    The central residue is determined by atom index 1: $N_i$ (residue $i$)

    \item \textbf{Psi ($\psi$)}: $[N_i, CA_i, C_i, N_{i+1}]$

    The central residue is determined by atom index 0: $N_i$ (residue $i$)
\end{itemize}

When MDTraj returns phi/psi arrays:
\begin{itemize}
    \item \texttt{phi\_idx}: array of shape $(n_{\phi}, 4)$ containing atom indices
    \item \texttt{phi}: array of shape $(n_{\text{frames}}, n_{\phi})$ containing angle values in radians
    \item Similarly for psi
\end{itemize}

The key insight: we must extract the \textit{residue index} from the atom quadruples to match phi and psi.

\subsection{Robust Alignment Algorithm}

\begin{algorithm}[H]
\caption{Robust $\phi/\psi$ Alignment}
\begin{algorithmic}[1]
\STATE Compute $\phi$ quadruples and values: \texttt{phi\_idx, phi = md.compute\_phi(traj)}
\STATE Compute $\psi$ quadruples and values: \texttt{psi\_idx, psi = md.compute\_psi(traj)}
\STATE Extract residue index for each $\phi$: \texttt{phi\_res[j] = residue of atom phi\_idx[j][1]}
\STATE Extract residue index for each $\psi$: \texttt{psi\_res[k] = residue of atom psi\_idx[k][0]}
\STATE Find common residues: $R_{\text{common}} = R_\phi \cap R_\psi$
\STATE Sort $R_{\text{common}}$ (to maintain sequence order)
\FOR{each residue $r \in R_{\text{common}}$}
    \STATE Find position in phi array: $j = \text{pos}_\phi(r)$
    \STATE Find position in psi array: $k = \text{pos}_\psi(r)$
    \STATE $\phi_{\text{aligned}}[:, i] = \phi[:, j]$
    \STATE $\psi_{\text{aligned}}[:, i] = \psi[:, k]$
\ENDFOR
\RETURN $\phi_{\text{aligned}}, \psi_{\text{aligned}}, R_{\text{common}}, \text{mask}$
\end{algorithmic}
\end{algorithm}

\subsection{Why This Matters}

Without proper alignment:
\begin{itemize}
    \item Token 0 might represent residue 2 for $\phi$ but residue 1 for $\psi$
    \item The model would try to learn relationships between \textit{unrelated} angles
    \item Training would fail or produce meaningless CVs
\end{itemize}

With alignment:
\begin{itemize}
    \item Token $i$ consistently represents residue $r_i$ for both angles
    \item The model learns true residue-level conformational patterns
    \item Attention weights correctly identify important residues
\end{itemize}

\subsection{Handling Edge Cases}

The alignment procedure naturally handles:
\begin{itemize}
    \item \textbf{Terminal residues}: First residue has no $\phi$, last has no $\psi$ $\rightarrow$ excluded from common set
    \item \textbf{Missing atoms}: Residues without backbone atoms (e.g., in broken chains) $\rightarrow$ excluded
    \item \textbf{Non-standard residues}: Handled transparently via topology information
    \item \textbf{Multi-chain proteins}: Each chain contributes its valid residues to the common set
\end{itemize}

The \texttt{mask} array (currently all True for valid residues) is kept for API consistency and future extensions (e.g., manually excluding flexible termini).

\subsection{Dataset and Temporal Split}

To avoid data leakage in MD trajectories (temporally correlated), a temporal split is used:
\begin{itemize}
    \item Training: first 90\% of frames
    \item Validation: last 10\% of frames
\end{itemize}

This ensures that the validation set represents system states never seen during training.

\section{Training: How the Model Learns}

\subsection{The Training Loop: Step by Step}

Training CVFormer involves iteratively updating model parameters to minimize the loss function. Here's what happens in each epoch:

\subsubsection{Forward Pass}
\begin{enumerate}
    \item \textbf{Batch loading}: Load a batch of $B$ conformations: $\mathbf{X} \in \mathbb{R}^{B \times N \times 4}$
    \item \textbf{Encoding}: Pass through encoder to get latent vectors: $\mathbf{z} \in \mathbb{R}^{B \times d_{\text{latent}}}$
    \item \textbf{Decoding}: Reconstruct dihedrals from latent: $\hat{\mathbf{X}} \in \mathbb{R}^{B \times N \times 4}$
    \item \textbf{Loss computation}: Calculate reconstruction error (and optionally contrastive loss)
\end{enumerate}

\subsubsection{Backward Pass}
\begin{enumerate}
    \item \textbf{Gradient computation}: Compute $\nabla_\theta \mathcal{L}$ via backpropagation
    \item \textbf{Gradient clipping}: Clip gradients to prevent exploding gradients (max norm = 1.0)
    \item \textbf{Parameter update}: Update weights using optimizer
    \item \textbf{Learning rate adjustment}: Update learning rate via scheduler
\end{enumerate}

\subsubsection{Validation and Early Stopping}
After each epoch:
\begin{enumerate}
    \item Evaluate on validation set (last 10\% of trajectory)
    \item Track validation loss, MAE$_\phi$, MAE$_\psi$
    \item If validation loss improves: save model weights, reset patience counter
    \item If validation loss doesn't improve: increment patience counter
    \item If patience counter reaches limit (e.g., 25 epochs): stop training
\end{enumerate}

This prevents overfitting and ensures the model generalizes to unseen conformations.

\subsection{Optimization Algorithm}

\subsubsection{AdamW Optimizer}
CVFormer uses AdamW (Adam with weight decay), which combines:
\begin{itemize}
    \item \textbf{Adaptive learning rates}: each parameter gets its own learning rate based on gradient history
    \item \textbf{Momentum}: smooths updates using exponential moving averages
    \item \textbf{Weight decay}: regularizes by penalizing large weights (prevents overfitting)
\end{itemize}

Update rule:
\begin{equation}
\theta_{t+1} = \theta_t - \eta_t \left(\frac{\hat{m}_t}{\sqrt{\hat{v}_t} + \epsilon} + \lambda \theta_t\right)
\end{equation}

where $\hat{m}_t$ is bias-corrected first moment, $\hat{v}_t$ is bias-corrected second moment, $\eta_t$ is learning rate, and $\lambda$ is weight decay coefficient (typically $10^{-4}$).

\subsubsection{Learning Rate Scheduler: Warmup + Cosine Decay}

A fixed learning rate often leads to unstable training (especially early on) or suboptimal convergence (late in training). CVFormer uses a two-phase scheduler:

\paragraph{Phase 1: Warmup (first 5 epochs)}
Linearly increase learning rate from 0 to $\eta_{\text{base}}$:
\begin{equation}
\eta_t = \eta_{\text{base}} \cdot \frac{t}{T_{\text{warmup}}}
\end{equation}

\textit{Why?} At initialization, weights are random. Large learning rates cause wild updates. Warmup allows the model to ``find its footing'' before aggressive optimization.

\paragraph{Phase 2: Cosine Annealing (remaining epochs)}
Gradually decrease learning rate following a cosine curve:
\begin{equation}
\eta_t = \eta_{\text{min}} + \frac{1}{2}(\eta_{\text{base}} - \eta_{\text{min}})\left(1 + \cos\left(\pi \frac{t - T_{\text{warmup}}}{T_{\text{total}} - T_{\text{warmup}}}\right)\right)
\end{equation}

where $\eta_{\text{min}} = 0.05 \cdot \eta_{\text{base}}$ (minimum learning rate).

\textit{Why cosine?} As training progresses, we want smaller updates to fine-tune the model. The smooth cosine decay:
\begin{itemize}
    \item Allows rapid learning early on (high $\eta$)
    \item Enables precise convergence later (low $\eta$)
    \item Avoids sudden drops that cause training instability
\end{itemize}

\begin{figure}[h]
\centering
\textit{[Learning rate schedule: linear warmup for 5 epochs, then cosine decay from $3 \times 10^{-4}$ to $1.5 \times 10^{-5}$ over remaining epochs]}
\end{figure}

\subsection{Recommended Hyperparameters}

\begin{table}[h]
\centering
\begin{tabular}{@{}lll@{}}
\toprule
\textbf{Parameter} & \textbf{Default} & \textbf{Description} \\
\midrule
\texttt{d\_model} & 64 & Embedding dimension \\
\texttt{nhead} & 4 & Number of attention heads \\
\texttt{num\_encoder\_layers} & 3 & Encoder layers \\
\texttt{num\_decoder\_layers} & 3 & Decoder layers \\
\texttt{dim\_feedforward} & 256 & Feedforward dimension \\
\texttt{dropout} & 0.2 & Dropout rate \\
\texttt{latent\_dim} & 2 & Latent space dimension \\
\texttt{memory\_tokens} & 4 & Memory tokens in decoder \\
\texttt{batch\_size} & 128 & Batch size \\
\texttt{lr} & 3e-4 & Learning rate (peak) \\
\texttt{warmup\_epochs} & 5 & Warmup epochs \\
\texttt{patience} & 25 & Early stopping patience \\
\bottomrule
\end{tabular}
\caption{Model and training hyperparameters}
\end{table}

\subsection{Training Procedure}

\begin{algorithm}[H]
\caption{Main Training Loop}
\begin{algorithmic}[1]
\STATE Load trajectory and compute aligned $\phi/\psi$
\STATE Create dataset with temporal split (90/10)
\STATE Initialize model, optimizer, scheduler
\FOR{epoch = 1 to max\_epochs}
    \FOR{batch in train\_loader}
        \STATE $\hat{\mathbf{X}}, \mathbf{z} = \text{model}(\mathbf{X}, \text{mask})$
        \STATE $\mathcal{L}_{\text{rec}} = \text{dihedral\_loss}(\hat{\mathbf{X}}, \mathbf{X}, \text{mask})$
        \IF{contrastive\_weight $> 0$}
            \STATE $\mathbf{X}' = \text{augment}(\mathbf{X})$
            \STATE $\mathbf{z}' = \text{model.encode}(\mathbf{X}', \text{mask})$
            \STATE $\mathcal{L}_c = \text{NT-Xent}(\mathbf{z}, \mathbf{z}')$
            \STATE $\mathcal{L} = \mathcal{L}_{\text{rec}} + \lambda_c \mathcal{L}_c$
        \ELSE
            \STATE $\mathcal{L} = \mathcal{L}_{\text{rec}}$
        \ENDIF
        \STATE Backpropagation and gradient clipping
        \STATE Update optimizer and scheduler
    \ENDFOR
    \STATE Evaluation on validation set
    \IF{validation loss improves}
        \STATE Save best model weights
        \STATE Reset patience counter
    \ELSE
        \STATE Increment patience counter
    \ENDIF
    \IF{patience counter $\geq$ max\_patience}
        \STATE Early stopping
    \ENDIF
\ENDFOR
\end{algorithmic}
\end{algorithm}


\section{Results Analysis: Interpreting the Learned CVs}

\subsection{What the Model Produces}

After training completes, CVFormer outputs several artifacts in the \texttt{output/} directory:

\begin{table}[h]
\centering
\small
\begin{tabular}{@{}ll@{}}
\toprule
\textbf{File} & \textbf{Content} \\
\midrule
\texttt{best\_model\_weights.pt} & Trained model parameters \\
\texttt{latents.npy} & CV values for all frames $(N_{\text{frames}} \times d_{\text{latent}})$ \\
\texttt{attn\_weights.npy} & Attention weights per frame $(N_{\text{frames}} \times N_{\text{residues}})$ \\
\texttt{attn\_weights\_mean.txt} & Mean attention weight per residue \\
\texttt{attn\_weights\_median.txt} & Median attention weight per residue \\
\texttt{token\_residue\_ids.txt} & Mapping from token index to residue ID \\
\texttt{attn\_importance\_by\_residue.txt} & Combined table of residue importance \\
\texttt{config.txt} & All hyperparameters used \\
\bottomrule
\end{tabular}
\caption{Output files from training}
\end{table}

\subsection{Understanding the Latent Space}

The \texttt{latents.npy} file contains the collective variables for each frame. For a 2D latent space:
\begin{itemize}
    \item Each row is a frame from the trajectory
    \item Column 0 is CV1, column 1 is CV2
    \item These are the coordinates you'll use for metadynamics biasing
\end{itemize}

\subsubsection{What Should You See?}

For a well-trained model on a protein with conformational transitions:

\paragraph{Good signs:}
\begin{itemize}
    \item \textbf{Clustered structure}: Different conformational states form distinct clusters in latent space
    \item \textbf{Smooth trajectories}: Transitions between states follow continuous paths (no random jumps)
    \item \textbf{Low reconstruction error}: MAE$_\phi$ and MAE$_\psi < 0.2$ radians ($\sim 11°$)
    \item \textbf{Clear separation}: If you know your protein has 2-3 states, you should see 2-3 clusters
\end{itemize}

\paragraph{Warning signs:}
\begin{itemize}
    \item \textbf{Uniform distribution}: CVs are random noise, not capturing conformational states
    \item \textbf{Collapsed latent space}: All points in a tiny region (model not using latent capacity)
    \item \textbf{High reconstruction error}: MAE $> 0.5$ radians (model can't compress effectively)
    \item \textbf{Validation loss diverging}: Model is overfitting to training data
\end{itemize}

If you see warning signs: try increasing \texttt{latent\_dim}, adding more encoder layers, training longer, or using contrastive loss.

\subsection{Analysis Notebooks}

\subsubsection{analyze-latent.ipynb}
Latent space analysis:
\begin{itemize}
    \item \textbf{Distributions}: 1D and 2D histograms of CVs
    \item \textbf{Clustering}: identification of discrete conformational states (K-means, DBSCAN)
    \item \textbf{Dynamics}: temporal evolution and state transitions
    \item \textbf{Correlations}: correlation analysis with other structural metrics (RMSD, Rg, etc.)
\end{itemize}

\subsubsection{residue-importance.ipynb}
Residue importance analysis:
\begin{itemize}
    \item \textbf{Global statistics}: mean, median, std of attention weights per residue
    \item \textbf{State-wise clustering}: identification of key residues for each conformational state
    \item \textbf{Transition analysis}: residues that change importance during transitions
    \item \textbf{Visualization}: heatmaps, ranked importance plots
\end{itemize}

\subsection{Attention Weight Interpretation: Finding Key Residues}

The attention pooling weights $\alpha_i$ are perhaps the most valuable output for biological insight. They answer: \textit{``Which residues are most important for defining conformational states?''}

\subsubsection{What Are These Weights?}

Recall from the encoder architecture:
\begin{equation}
\mathbf{z}_{\text{pool}} = \sum_{i=1}^{N} \alpha_i \mathbf{h}_i
\end{equation}

The weight $\alpha_i$ tells us how much residue $i$ contributes to the pooled representation (and thus to the final CV values). Higher $\alpha_i$ means residue $i$ has more influence.

\subsubsection{Two Levels of Analysis}

\paragraph{1. Global Importance (Mean/Median Weights)}
The files \texttt{attn\_weights\_mean.txt} and \texttt{attn\_weights\_median.txt} contain average weights across all frames. These identify residues that are \textit{consistently important} throughout the simulation.

\textbf{Biological interpretation:}
\begin{itemize}
    \item High mean weight: residue participates in major conformational changes
    \item Examples: hinge residues in domain motions, active site in binding events, allosteric coupling sites
\end{itemize}

\paragraph{2. State-Specific Importance (Per-Frame Weights)}
The file \texttt{attn\_weights.npy} contains weights for each frame. By clustering frames in latent space and averaging weights per cluster, you can identify:
\begin{itemize}
    \item Residues important for \textit{specific} conformational states
    \item Residues that change importance during transitions
    \item Different sets of residues controlling different motions
\end{itemize}

\subsubsection{Practical Use Cases}

\paragraph{Drug Design}
High-attention residues near a binding pocket indicate where conformational flexibility matters most. These are candidates for:
\begin{itemize}
    \item Allosteric drug targeting
    \item Mutagenesis studies
    \item Understanding drug resistance mechanisms
\end{itemize}

\paragraph{Protein Engineering}
If you want to stabilize a particular conformation:
\begin{itemize}
    \item Identify high-attention residues
    \item Mutate them to favor desired $\phi/\psi$ angles
    \item Or introduce cross-links to restrict motion
\end{itemize}

\paragraph{Mechanistic Understanding}
Compare attention weights with known structural features:
\begin{itemize}
    \item Do high-attention residues coincide with experimental data (e.g., NMR dynamics)?
    \item Are they in expected regions (domain interfaces, loops)?
    \item Do they reveal unexpected coupling mechanisms?
\end{itemize}

\subsubsection{Example Interpretation}

Suppose CVFormer identifies residues 45, 67, and 89 as having highest mean attention weights:
\begin{enumerate}
    \item \textbf{Check structure}: Locate these residues in 3D structure
    \item \textbf{Analyze dynamics}: Compute RMSD/RMSF for these residues across trajectory
    \item \textbf{Correlate with CVs}: Plot $\phi_{45}$, $\psi_{45}$ vs CV1 — is there strong correlation?
    \item \textbf{Biological context}: Are these in functional regions (active site, binding interface, etc.)?
\end{enumerate}

The attention mechanism essentially performs an automated version of what an expert would do manually: identify the few residues that ``matter'' for conformational behavior.

\section{Practical Usage}

\subsection{Installation}

\begin{lstlisting}[language=bash]
# Create conda environment
mamba env create -f environment.yml
mamba activate cvformer
\end{lstlisting}

The environment includes:
\begin{itemize}
    \item Python 3.10
    \item PyTorch (CPU or CUDA)
    \item MDTraj for preprocessing
    \item NumPy, Matplotlib for analysis
\end{itemize}

\subsection{Basic Training}

\begin{lstlisting}[language=bash]
python testit.py \
  --trajectory data/traj_skip100NoH.xtc \
  --topology data/2JOF-0-proteinNoH.pdb \
  --output_dir output \
  --latent_dim 2 \
  --epochs 1000 \
  --batch_size 128 \
  --patience 25
\end{lstlisting}

\subsection{Training with Contrastive Learning}

\begin{lstlisting}[language=bash]
python testit.py \
  --trajectory data/traj_skip100NoH.xtc \
  --topology data/2JOF-0-proteinNoH.pdb \
  --output_dir output_contrastive \
  --latent_dim 2 \
  --contrastive_weight 0.1 \
  --contrastive_temp 0.1 \
  --contrastive_noise 0.05
\end{lstlisting}


\section{Implementation Improvements}

\subsection{Evolution from Main Branch}

The \texttt{guglielmo\_contrastive\_implementation} branch introduces significant improvements over \texttt{main}:

\begin{table}[h]
\centering
\small
\begin{tabular}{@{}lll@{}}
\toprule
\textbf{Aspect} & \textbf{Main} & \textbf{Guglielmo Implementation} \\
\midrule
Structure & Monolithic (single file) & Modular (pkgs/) \\
Dihedral alignment & Simple & Robust with mask \\
Normalization & Basic & Unit circle projection \\
Scheduler & Fixed & Warmup + Cosine decay \\
Metrics & MSE & MSE + Angular MAE \\
Interpretability & - & Attention weights export \\
Analysis & - & Notebook suite \\
Contrastive loss & - & SimCLR-style optional \\
\bottomrule
\end{tabular}
\caption{Implementation comparison}
\end{table}

\subsection{Key Advantages}

\begin{enumerate}
    \item \textbf{Robustness}: automatic handling of complex topologies and termini
    \item \textbf{Interpretability}: extraction and analysis of attention weights
    \item \textbf{Stability}: unit circle normalization and gradient clipping
    \item \textbf{Flexibility}: modular architecture for easy extensions
\end{enumerate}

\section{Limitations and Future Developments}

\subsection{Current Limitations}
\begin{itemize}
    \item The model is purely deterministic (no variational autoencoder)
    \item Requires sufficiently long training trajectories (thousands of frames)
    \item Attention pooling uses soft weights (not sparse)
    \item Does not directly support multi-chain or heterogeneous systems
\end{itemize}

\subsection{Possible Extensions}
\begin{itemize}
    \item \textbf{VAE variant}: variational implementation for sampling and uncertainty quantification
    \item \textbf{Hierarchical models}: multi-scale encoder to capture dependencies at different levels
    \item \textbf{Graph Neural Networks}: integration of topological information (contacts, distances)
    \item \textbf{Multi-task learning}: joint training on different properties (stability, function, etc.)
    \item \textbf{Active learning}: automatic selection of configurations for enhanced sampling
\end{itemize}

\section{Conclusions}

\subsection{Summary: The Complete CVFormer Pipeline}

CVFormer provides an end-to-end solution for automatic CV discovery from MD trajectories. The complete workflow is:

\begin{enumerate}
    \item \textbf{Input}: MD trajectory (any format supported by MDTraj) + topology file
    \item \textbf{Preprocessing}: Robust alignment of $\phi/\psi$ angles, conversion to $(\sin, \cos)$ representation
    \item \textbf{Training}: Autoencoder learns to compress conformational information into 2-3 latent dimensions
    \item \textbf{Analysis}: Extract CVs and attention weights, visualize conformational landscape
    \item \textbf{Interpretation}: Identify key residues driving conformational changes
    \item \textbf{Export}: CVs can be saved and used for further analysis or custom biasing implementations
\end{enumerate}

\subsection{Key Innovations}

CVFormer advances beyond traditional autoencoder approaches through:

\begin{itemize}
    \item \textbf{Transformer architecture}: Captures long-range residue dependencies via self-attention
    \item \textbf{Attention pooling}: Provides interpretable importance scores for each residue
    \item \textbf{Proper periodicity handling}: $(\sin, \cos)$ representation + unit circle normalization
    \item \textbf{Robust preprocessing}: Automatic handling of terminal residues, missing atoms, multi-chain systems
    \item \textbf{Contrastive learning}: Optional robustness enhancement via SimCLR-style loss
    \item \textbf{Comprehensive output}: Latent variables, attention weights, and detailed metrics
\end{itemize}

\subsection{When to Use CVFormer}

CVFormer is most effective when:

\paragraph{Good fit:}
\begin{itemize}
    \item You have a trajectory (5000+ frames) covering multiple conformational states
    \item The conformational dynamics are driven by backbone dihedrals (not just sidechain or overall shape)
    \item You want interpretable CVs (which residues matter)
    \item You need to identify key residues for conformational transitions
\end{itemize}

\paragraph{Consider alternatives if:}
\begin{itemize}
    \item Conformational changes are primarily sidechain-driven (use tools that include $\chi$ angles)
    \item You have very limited data (<1000 frames)
    \item The system is highly disordered (IDPs may need different representations)
    \item You only need dimensionality reduction for visualization (PCA/t-SNE may suffice)
\end{itemize}

\subsection{Main Strengths}

\begin{enumerate}
    \item \textbf{Automation}: No manual CV selection required
    \item \textbf{Interpretability}: Attention weights reveal \textit{which} residues matter
    \item \textbf{Robustness}: Handles complex topologies, multi-chain systems, missing residues
    \item \textbf{Extensibility}: Modular design supports easy modifications (new loss functions, architectures, etc.)
\end{enumerate}

\subsection{Looking Forward}

CVFormer represents a step toward \textit{automated} conformational analysis, where:
\begin{itemize}
    \item CVs are discovered from data rather than expert intuition
    \item The model tells you which parts of the protein are important
    \item The learned representation provides interpretable insights into protein dynamics
\end{itemize}

The approach is general and can be extended to:
\begin{itemize}
    \item Include sidechain dihedrals ($\chi$ angles)
    \item Incorporate distance-based features (contacts, hydrogen bonds)
    \item Multi-task learning (predict binding affinity, stability, etc.)
    \item Transfer learning (pre-train on many proteins, fine-tune on target)
\end{itemize}

The project is continuously evolving and open to contributions to improve robustness, performance, and applicability to increasingly complex biomolecular systems.

\section*{References}

\begin{itemize}
    \item Vaswani et al., ``Attention Is All You Need'', NeurIPS 2017
    \item Chen et al., ``A Simple Framework for Contrastive Learning of Visual Representations'' (SimCLR), ICML 2020
    \item Bonomi et al., ``PLUMED: A portable plugin for free-energy calculations with molecular dynamics'', Computer Physics Communications 2009
    \item McGibbon et al., ``MDTraj: A Modern Open Library for the Analysis of Molecular Dynamics Trajectories'', Biophysical Journal 2015
\end{itemize}

\end{document}
